\scanpage{138}\mspage{- 137 -}\hfill IX

sur $K$, qui jouera dans le cas présent le même rôle que 11.2 dans le cas traité précédemment :

\uline{Théorème} 11.5. \uline{Supposons} $S$ \uline{hensélien}, \uline{et que} $A_K$ \uline{ait réduction semi-stable sur} $S$. \uline{Soit} $\ell$ \uline{un nombre premier}, \uline{et considérons l'accouplement de monodromie} (9.1.2), \uline{d'où un homomorphisme canonique}
\[
 \widetilde u_\ell:\underline M_\ell\longrightarrow\underline{\check M}'_\ell\quad .
\]
\uline{Soit} $\Phi(\ell)=\mathrm{coker}\,\widetilde u_\ell$ ; \uline{alors on a un isomorphisme canonique}
\[
 \Phi_{\unsure{o|0}}(\ell)\simeq\Phi(\ell)\times_S s\quad .\tag*{(11.5.1)}
\]

De cet énoncé résulte immédiatement la définition d'un accouplement parfait (11.4.1), puisque on trouve, en appliquant 11.5 à $A'_K$ au lieu de $A_K$ et utilisant (10.2.2), qu'on a aussi un isomorphisme canonique
\[
 \Phi'_o(\ell)\simeq\Phi'(\ell)\times_S s\quad,\quad\text{où }\Phi'(\ell)=\mathrm{coker}({}^t\widetilde u_\ell:\underline M'_\ell\longrightarrow\underline{\check M}_\ell)\quad .
\]

\uline{Remarques} 11.5.2. a) Le théorème 11.5 implique notamment que $\widetilde u_\ell$ est une isogénie i.e. que $u_\ell$ est séparant, y compris dans le cas $\ell=p$, où cette assertion n'est pas triviale sur la définition de 9.6. On notera que nous n'\deleted{\illegible{1}}utilisons pas ici le théorème (nettement plus profond) 10.4, qui sera prouvé au par. suivant.

b)\quad Utilisant 10.4, qui nous fournit un accouplement de monodromie $u:\underline M\otimes\underline M'\longrightarrow\mathbb Z_S$, et désignant par $\Phi$ le $\underleftarrow{\text{\deleted{\unsure{o|c|g}}\handwritten{\inserted{\unsure{c}}}onoyau}}$ de l'homomorphisme correspondant
\[
 \widetilde u:\underline M\longrightarrow\underline{\check M}'\quad ,
\]
\margin{\unsure{2}}
% archive sheet number: 137
% stamp: numbered label 455.
